\title{DeepSeekMath: Pushing the Limits of Mathematical Reasoning in Open Language Models}

\begin{abstract}
The intricacies and structured demands of mathematical reasoning present considerable obstacles for language models. In this study, we present DeepSeekMath~7B, an extension of DeepSeek-Coder-Base-v1.5 7B that undergoes continued pre-training with 120B mathematics-focused tokens extracted from Common Crawl, in addition to natural language and programming code datasets. DeepSeekMath~7B attains a notable score of 51.7\% on the rigorous MATH competition benchmark, accomplishing this without the use of supplementary toolkits or ensemble methods, and nearing the benchmarks set by Gemini-Ultra and GPT-4. When employing self-consistency over 64 samples, DeepSeekMath~7B further elevates performance to 60.9\% on the MATH dataset. The notable proficiency of DeepSeekMath~in mathematical reasoning is underpinned by two principal strategies: Firstly, we fully exploit the wealth of mathematics-relevant web data via a carefully designed data selection process. Secondly, we propose Group Relative Policy Optimization (GRPO), an adaptation of Proximal Policy Optimization (PPO), to reinforce mathematical reasoning capacities while simultaneously improving PPO’s memory efficiency.
\end{abstract}

\section{Introduction}

The advent of large language models (LLMs) has dramatically transformed methodologies for tackling mathematical reasoning within artificial intelligence, leading to remarkable progress on benchmarks assessing quantitative reasoning \citep{MATH} as well as geometric reasoning \citep{trinh2024solving}. Additionally, these LLMs have served as valuable tools in supporting human efforts to address intricate mathematical tasks \citep{tao}. Nonetheless, advanced models such as GPT-4 \citep{gpt4} and Gemini-Ultra \citep{gemini} are not freely available to the public, and open-source alternatives presently lag significantly in their comparative performance.

This paper presents DeepSeekMath, a specialized language model tailored for mathematics that surpasses other open-source models in mathematical tasks and nearly matches GPT-4’s performance on established academic benchmarks. Central to this achievement is the construction of the DeepSeekMath Corpus, an extensive and carefully curated pre-training dataset containing 120B math tokens. The corpus is assembled from the Common Crawl (CC) through the use of a fastText-based classifier \citep{joulin2016fasttext}. Initially, the classifier is trained on positive samples from OpenWebMath \citep{openwebmath}, alongside a wide variety of web pages designated as negative samples. With this trained model, we identify additional math-related content within the CC, which is subsequently verified and improved through manual annotation. The classifier undergoes retraining with this expanded and higher quality data to further enhance its effectiveness. Our empirical results demonstrate the superior quality of the collected corpus: DeepSeekMath-Base 7B attains 64.2\% accuracy on GSM8K \citep{gsm8k} and 36.2\% on the advanced MATH competition dataset \citep{MATH}, both scores surpassing those achieved by Minerva 540B \citep{minerva}. Furthermore, thanks to the multilingual design of the DeepSeekMath Corpus, we observe notable gains in Chinese mathematical benchmarks \citep{wei2023cmath,agieval}. We propose that our methods for assembling and refining mathematical datasets can serve as a foundation for subsequent research, while recognizing that considerable potential for future advancements remains.

DeepSeekMath-Base is built upon DeepSeek-Coder-Base-v1.5 7B \citep{deepseek-coder}, as empirical evidence suggests that initializing from a code-specialized model offers a stronger foundation than a conventional language model. We also find that mathematical pre-training enhances not only mathematical problem-solving but also boosts results on MMLU \citep{mmlu} and BBH \citep{bbh} evaluations, thereby strengthening both mathematical proficiency and general reasoning skills.

Following pre-training, we conduct mathematical instruction tuning on DeepSeekMath-Base using data involving chain-of-thought \citep{cot}, program-of-thought \citep{pot,pal}, and tool-integrated reasoning \citep{tora} approaches. The resulting model, DeepSeekMath-Instruct 7B, outperforms all other models in the 7B parameter class and achieves performance on par with the best available open-source instruction-tuned models with 70B parameters.

Additionally, we propose Group Relative Policy Optimization (GRPO), an RL variant based on Proximal Policy Optimization (PPO) \citep{schulman2017proximal}. Unlike standard PPO, GRPO does not utilize a critic network; instead, it estimates baselines using aggregated group scores, significantly cutting computational costs. With only a selected subset of English instruction data, GRPO yields notable improvements over the already strong DeepSeekMath-Instruct, demonstrated by gains on both in-domain (GSM8K: 82.9\% $\rightarrow$ 88.2\%, MATH: 46.8\% $\rightarrow$ 51.7\%) and out-of-domain tasks (such as CMATH: 84.6\% $\rightarrow$ 88.8\%) during RL training. Moreover, we offer a unifying framework to interpret approaches like Rejection Sampling Fine-Tuning (RFT) \citep{yuan2023scaling}, Direct Preference Optimization (DPO) \citep{dpo}, PPO, and GRPO, classifying these as forms of direct or simplified RL. We further undertake extensive empirical studies—contrasting, for instance, online versus offline updates, supervision by outcome versus process, and single-turn against iterative RL—to clarify the core components underlying this paradigm. Finally, we elucidate the mechanisms through which RL improves instruction-tuned models, and outline future research directions for achieving more effective RL within this unified framework.

\subsection{Contributions}
This work advances the field through both large-scale mathematical pre-training and an in-depth investigation of reinforcement learning techniques.

\textbf{Math Pre-Training at Scale}

\begin{itemize}[topsep=0pt]
    \item We demonstrate that the Common Crawl dataset, available publicly, is a rich resource for mathematical data. By employing a rigorous and thorough data filtering approach, we create the DeepSeekMath Corpus: a curated dataset containing 120B tokens derived from math-focused web content. This corpus is approximately seven times larger than that used in Minerva \citep{minerva}, and nine times the volume of the newly published OpenWebMath \citep{openwebmath}.

    \item DeepSeekMath-Base 7B, our base pre-trained model, matches the performance of Minerva 540B \citep{minerva}. This result implies that sheer parameter count does not solely determine a model’s proficiency in mathematical reasoning; models of smaller scale, if trained on highly relevant data, can deliver equally competitive outcomes.

    \item We present insights gained from our math model training processes. Pre-training on code before commencing math-specific training significantly enhances a model’s mathematical problem-solving abilities, regardless of whether external tools are involved. This finding provides a partial resolution to the open question: \textit{does code pre-training contribute to reasoning skills?} Our evidence suggests a positive effect, at least for mathematics.

    \item Contrary to prevailing trends, we observe that incorporating arXiv paper data—despite its frequent use in mathematical research—fails to yield significant improvements across all the mathematical benchmarks tested in this study.
\end{itemize}

\textbf{Exploration and Analysis of Reinforcement Learning}

\begin{itemize}[topsep=0pt]
    \item We present Group Relative Policy Optimization (GRPO), a novel reinforcement learning approach that achieves both efficiency and efficacy. Unlike conventional methods that rely on a critic network, GRPO leverages group-based score estimations for baseline calculation, substantially decreasing the computational demands compared to Proximal Policy Optimization (PPO).
    \item Our empirical findings reveal that applying GRPO considerably improves the outcomes of our instruction-tuned system, DeepSeekMath-Instruct, utilizing exclusively instruction-tuning datasets. Additionally, we notice increased robustness and transferability, as out-of-domain performance also improves during the reinforcement learning phase.
    \item We offer an integrated theoretical framework that elucidates the connections among various methodologies, including RFT, DPO, PPO, and GRPO. Through comprehensive experimentation—spanning online versus offline optimization, outcome versus process-based supervision, and single-turn versus iterative learning—we rigorously analyze the fundamental aspects that shape this unified framework.
    \item Building upon this integrative perspective, we investigate the underlying factors contributing to the success of reinforcement learning and propose several promising avenues to further advance reinforcement learning methods for large language models (LLMs).
\end{itemize}

\subsection{Summary of Evaluations and Metrics}

\begin{itemize}[topsep=0pt]
    \item \textbf{English and Chinese Mathematical Reasoning}: We systematically evaluate our models on a diverse set of mathematical reasoning datasets in both English and Chinese, spanning problem sets from elementary to collegiate levels. The English benchmarks utilized include GSM8K \citep{gsm8k}, MATH \citep{MATH}, SAT \citep{llemma}, OCW Courses \citep{minerva}, and MMLU-STEM \citep{mmlu}. For Chinese, we incorporate MGSM-zh \citep{mgsm}, CMATH \citep{wei2023cmath}, Gaokao-MathCloze \citep{agieval}, and Gaokao-MathQA \citep{agieval}. Our evaluation framework assesses the model's competence in independently generating step-by-step solutions in natural language without external tools, as well as its ability to resolve mathematical tasks programmatically using Python.

    On the English evaluation sets, DeepSeekMath-Base achieves results on par with Minerva 540B \citep{minerva}, a prominent proprietary model, and markedly outperforms all other open-source foundational models, such as Mistral 7B \citep{mistral} and Llemma-34B \citep{llemma}, regardless of prior mathematical pre-training. In particular, the model demonstrates a notable advantage on Chinese datasets, likely attributed to our strategy of integrating diverse, high-quality multilingual math data rather than limiting pre-training to English sources, in contrast with preceding studies \citep{minerva,llemma}. Furthermore, after applying mathematical instruction tuning and reinforcement learning, DeepSeekMath-Instruct and DeepSeekMath-RL achieve outstanding results, surpassing the 50\% accuracy threshold on the challenging MATH dataset for the first time among open-source models.
\end{itemize}

\item \textbf{Formal Mathematics}: We assess the capabilities of DeepSeekMath-Base on informal-to-formal theorem proving using the benchmark introduced by \citep{dsp_proof}, focusing on the miniF2F dataset \citep{minif2f} and utilizing Isabelle \citep{isabelle} as the automated proof assistant. Results indicate that DeepSeekMath-Base achieves impressive few-shot performance in the task of autoformalization.
\item \textbf{Natural Language Understanding, Reasoning, and Code}: To comprehensively evaluate the model’s proficiency in language comprehension, reasoning, and code generation, we conduct assessments on the Massive Multitask Language Understanding (MMLU) benchmark \citep{mmlu}—which comprises 57 varied multiple-choice tasks—along with BIG-Bench Hard (BBH) \citep{bbh}, consisting of 23 tasks primarily emphasizing complex, multi-step reasoning. Additionally, we evaluate with HumanEval \citep{codex} and MBPP \citep{mbpp}, two widely adopted benchmarks for code-oriented language models. Math-centric pre-training is observed to contribute positively to both linguistic understanding and reasoning performance.

\section{Math Pre-Training}

\subsection{Data Collection and Decontamination}
This section details the methodology used to build the DeepSeekMath Corpus by leveraging Common Crawl data. As illustrated in Figure \ref{fig:data_collect}, we employ an iterative pipeline, systematically assembling a large-scale corpus with mathematical content. The process begins with a high-quality seed dataset, and this general methodology can readily be adapted to domains such as programming.

Our starting point is OpenWebMath \citep{openwebmath}, a carefully curated collection of mathematical texts from the web, which serves as our initial seed corpus. Using this resource, we train a fastText model \citep{joulin2016fasttext} to retrieve additional mathematical web content analogous to OpenWebMath. For training, we randomly sample 500,000 examples from the seed dataset as positive cases, and another 500,000 web pages from Common Crawl to serve as negative cases. Training is conducted with an open-source library\footnote{\url{https://fasttext.cc}}, configuring the model with a vector size of 256, a learning rate of 0.1, a word n-gram maximum length of 3, a minimum word frequency of 3, and running for 3 epochs. To curtail the Common Crawl’s volume, we apply URL-based deduplication and near-deduplication, bringing the dataset down to 40B HTML pages. Utilizing the fastText model, we identify additional math-related web pages from this deduplicated subset. To filter for high-quality mathematical material, the retrieved pages are ranked according to the fastText scores, with only the highest scoring ones retained. Pre-training trials are performed on the top 40B, 80B, 120B, and 160B tokens to estimate the appropriate data volume. For the initial cycle, the top 40B tokens are selected for preservation.

Following the initial round of data collection, a significant number of mathematical web pages were still missing. This was primarily due to the limited diversity of positive examples used to train the fastText model. To address this, we sought out additional mathematical web sources to enhance the initial seed dataset and thereby improve the fastText model’s performance. Concretely, we partitioned the entire Common Crawl dataset into distinct domains, where each domain corresponds to web pages under the same base URL. For each domain, we assessed the proportion of pages acquired during the first collection round. Domains with more than 10\% of pages captured (such as \url{mathoverflow.net}) were labeled as math-oriented.

Next, we performed manual annotation of URL paths within these math-related domains that point specifically to mathematical content (e.g., \url{mathoverflow.net/questions}). Web pages corresponding to these identified URLs, but not yet included in the dataset, were then incorporated into our seed corpus. By augmenting the pool of positive examples in this way, we were able to further train the fastText model for better recall in subsequent data collection rounds. Ultimately, after four rounds, we obtained a total of 35.5M mathematical web pages, containing approximately 120B tokens. By the fourth round, we observed that nearly 98\% of the total data had already been obtained by the end of the third round, leading us to terminate further collection.

In order to prevent any overlap with benchmark datasets, we adopt the filtering procedure from \cite{deepseek-coder}, eliminating any web pages containing question or answer material from English mathematical benchmarks (such as GSM8K \citep{gsm8k} and MATH \citep{MATH}) as well as Chinese benchmarks (including CMATH \citep{wei2023cmath} and AGIEval \citep{agieval}). Our filtering process is as follows: if any sequence of ten consecutive tokens (10-grams) in a text segment matches exactly with any substring from the benchmarks, that segment is excluded from the math training data. For shorter benchmark samples with lengths between three and nine tokens, we apply the same exact matching criterion to remove any contaminated web pages.

\subsection{Assessing the DeepSeekMath~Corpus Quality}

We conduct pre-training studies to compare the DeepSeekMath~Corpus with other recently introduced math-specific corpora:

\begin{itemize}[topsep=0pt]
    \item \textbf{MathPile} \citep{mathpile}: This composite dataset (8.9B tokens) aggregates content from textbooks, Wikipedia, ProofWiki, CommonCrawl, StackExchange, and arXiv, with arXiv accounting for over 85\% of its content.
    \item \textbf{OpenWebMath} \citep{openwebmath}: Comprises 13.6B tokens derived from mathematical content extracted from CommonCrawl.
    \item \textbf{Proof-Pile-2} \citep{llemma}: A collection including OpenWebMath, AlgebraicStack (10.3B tokens of mathematical code), and arXiv articles (28.0B tokens). For Proof-Pile-2 experiments, we adhere to the arXiv:Web:Code ratio of 2:4:1 as recommended by \cite{llemma}.
\end{itemize}

\subsubsection{Training Setting}
\label{sec:quality-policy}
We conduct mathematics-specific training on a broadly pre-trained language model comprising 1.3B parameters, aligned in architecture with the DeepSeek LLMs \citep{deepseek-llm}—hereafter referred to as DeepSeek-LLM 1.3B. For each mathematical dataset, a dedicated model is trained across 150B tokens. All training runs utilize the lightweight and efficient HAI-LLM framework \citep{haillm}. Consistent with the protocol for DeepSeek LLMs, we adopt the AdamW optimizer \citep{adamW} with parameter settings $\beta_1=0.9$, $\beta_2=0.95$, and $\mathrm{weight\_decay}=0.1$. A multi-phase learning rate schedule is employed: the rate peaks after a warmup of 2,000 steps, drops to 31.6\% of the maximum after 80\% of the training progress, and subsequently falls to 10.0\% of its peak value by the 90\% completion mark. The peak learning rate is capped at 5.3e-4. For all experiments, we use a batch size of 4 million tokens and a context window of 4,000 tokens.

\subsubsection{Evaluation Results}

\textbf{The DeepSeekMath~Corpus is distinguished by its high quality, extensive multilingual content, and superior scale.}

\begin{itemize}[topsep=0pt]
    \item \textbf{High-quality}:
    We assess model performance across eight mathematical benchmarks using few-shot chain-of-thought prompting \cite{cot}. As demonstrated in Table \ref{tab:corpora_comparison}, the model trained on DeepSeekMath~Corpus achieves consistently superior results compared to alternatives. In Figure \ref{fig:corpora_comparisons}, performance for DeepSeekMath-trained models surpasses that of those trained on Proof-Pile-2 at the 50B token mark (equivalent to one full pass over Proof-Pile-2), suggesting that the average sample quality in DeepSeekMath is notably higher.
    \item \textbf{Multilingual}:
    The DeepSeekMath~Corpus contains data in a variety of languages, with English and Chinese comprising the bulk of entries. According to Table \ref{tab:corpora_comparison}, models trained on DeepSeekMath~Corpus demonstrate significant gains in mathematical reasoning for both English and Chinese, unlike previous mathematical corpora—predominantly focused on English—which yield minimal improvement or even deteriorated performance for Chinese-language tasks.
    \item \textbf{Large-scale}:
    DeepSeekMath~Corpus eclipses existing mathematical corpora by several multiples in data size. Figure \ref{fig:corpora_comparisons} illustrates that DeepSeek-LLM 1.3B models trained on DeepSeekMath exhibit more rapid progress and sustained improvement over time. Conversely, smaller baseline corpora have undergone multiple training passes, resulting in early saturation of model performance.
\end{itemize}

\subsection{Training and Evaluating DeepSeekMath-Base 7B}

This section details the development of DeepSeekMath-Base 7B, a foundational model distinguished by its advanced mathematical reasoning skills. The model is initialized with DeepSeek-Coder-Base-v1.5 7B \citep{deepseek-coder} and further trained on a corpus totaling 500B tokens. The training data is sourced as follows: 56\% from the DeepSeekMath Corpus, 4\% from AlgebraicStack, 10\% from arXiv, 20\% from Github code repositories, and the final 10\% from Common Crawl-derived natural language text in both English and Chinese. Training procedures generally follow those outlined in Section \ref{sec:quality-policy}, with two main modifications: the maximum learning rate is increased to 4.2e-4, and each batch processes 10M tokens.

We implement an extensive evaluation of DeepSeekMath-Base 7B’s mathematical competencies. This includes examining its proficiency in independently generating comprehensive mathematical solutions, applying tool-assisted problem solving, and performing formal theorem proving. To establish a broader understanding of the model’s utility, we also assess its general capabilities, including natural language comprehension, logical reasoning, and code generation.

\paragraph{Mathematical Problem Solving with Step-by-Step Reasoning}

To assess the model’s approach to mathematical problem solving, we employ few-shot chain-of-thought prompting \citep{cot} over eight diverse English and Chinese benchmarks. These benchmarks are selected to test quantitative reasoning (such as GSM8K \citep{gsm8k}, MATH \citep{MATH}, and CMATH \citep{wei2023cmath}) and multiple-choice tasks (for example, MMLU-STEM \citep{mmlu} and Gaokao-MathQA \citep{agieval}), thereby spanning a range of mathematical subjects and problem difficulty levels, from foundational to collegiate mathematics.

Table \ref{tab:base_math_cot} presents results showing that DeepSeekMath-Base 7B attains the highest scores across all eight evaluated benchmarks compared to other open-source base models, such as the broadly-adopted Mistral 7B \citep{mistral} and the recently introduced Llemma 34B \citep{llemma}, which was trained with Proof-Pile-2 \citep{llemma}. Particularly noteworthy is its performance on the MATH competition dataset, where DeepSeekMath-Base exceeds other open-source base models by a margin of more than 10\% in absolute terms, and even outperforms the proprietary Minerva 540B \citep{minerva}—a base model 77 times larger that builds upon PaLM \citep{palm} with additional mathematics-specific training.

\paragraph{Mathematical Problem Solving with Tool Use} For assessments involving program-assisted mathematical reasoning, we conduct experiments on GSM8K and MATH employing few-shot program-of-thought prompting \citep{pot,pal}. The models are instructed to address each problem by crafting a Python script, making use of libraries like \textit{math} and \textit{sympy} for advanced calculations. The solution is derived by running the program and evaluating its output as the answer. Table \ref{tab:base_math_pot_proof} reveals that DeepSeekMath-Base 7B achieves superior results compared to the prior best-performing open-source model, Llemma 34B.

\paragraph{Formal Mathematics}

The automation of formal proof generation serves to improve the accuracy, dependability, and efficiency of mathematical reasoning, garnering growing attention lately. We assess the capability of DeepSeekMath-Base 7B in the informal-to-formal proof task outlined in \citep{dsp_proof}, where the model generates formal proofs given an informal statement, its formal version, and an accompanying informal proof. Our evaluation uses miniF2F \citep{minif2f}, a standard for Olympiad-level formal mathematics, requiring the production of Isabelle-formalized proofs for each item via few-shot prompting. Following \cite{dsp_proof}, the models generate outlines (proof sketches) which are then completed with the automated prover Sledgehammer \citep{sledgehammer}. Results reported in Table \ref{tab:base_math_pot_proof} highlight DeepSeekMath-Base 7B’s strong capabilities in automated proof formalization.

\paragraph{Natural Language Understanding, Reasoning, and Code}

To measure performance in natural language comprehension, reasoning, and programming, we use MMLU \citep{mmlu} for language understanding, BBH \citep{bbh} for reasoning, and HumanEval \citep{codex} and MBPP \citep{mbpp} for code generation. As reported in Table \ref{tab:understanding_reasoning_code}, DeepSeekMath-Base 7B shows marked improvement over its predecessor, DeepSeek-Coder-Base-v1.5 \citep{deepseek-coder}, in both MMLU and BBH tasks, illustrating the positive effects of specialized mathematical training on both language comprehension and reasoning ability. The incorporation of code-specific tokens in continual training enables DeepSeekMath-Base 7B to preserve the coding performance level of DeepSeek-Coder-Base-v1.5 across both coding benchmarks. In aggregate, DeepSeekMath-Base 7B significantly outperforms the general-purpose Mistral 7B \citep{mistral} in all three benchmarks targeting reasoning and programming abilities.

\section{Supervised Fine-Tuning}

\subsection{SFT Data Curation}

To support instruction tuning, we assemble a mathematical dataset that encompasses both English and Chinese tasks drawn from multiple mathematical disciplines and spanning a wide spectrum of difficulty. Each problem is matched with an explanation formatted in chain-of-thought (CoT) \citep{cot}, program-of-thought (PoT) \citep{pot,pal}, or with step-by-step tool-enhanced reasoning \citep{tora}. The combined dataset contains a total of 776K training samples.

\begin{itemize}[topsep=0pt]
    \item \textbf{English mathematical datasets}: Tool-integrated answer annotations are provided for GSM8K and MATH tasks. In addition, a portion of MathInstruct \citep{MathInstruct} and the Lila-OOD \citep{lila} training set are included, both featuring solutions constructed with CoT or PoT methods. Our English-language dataset is comprehensive, featuring a variety of mathematical topics, such as algebra, probability, number theory, calculus, and geometry.
    \item \textbf{Chinese mathematical datasets}: We curate Chinese-language mathematical problems for K-12 students, distributed over 76 sub-domains like linear equations. These are accompanied by answers formatted as both CoT and tool-based reasoning.
\end{itemize}

\subsection{Training and Evaluating DeepSeekMath-Instruct 7B}

Here, we present DeepSeekMath-Instruct 7B, a model refined for mathematical instruction following its pretraining as DeepSeekMath-Base. To create the input, we concatenate randomly selected training instances up to a limit of 4K tokens. Model training is conducted for 500 steps, utilizing a batch size of 256 and maintaining a fixed learning rate of 5e-5.

For evaluation, we assess the model's mathematical reasoning abilities with and without tool access on four quantitative reasoning benchmarks in both English and Chinese. The results are compared against top-performing models contemporaneous to our study:

\begin{itemize}[topsep=0pt]
    \item \textbf{Closed-source models} comprise: (1) the GPT suite, highlighting GPT-4 \citep{gpt4} and GPT-4 Code Interpreter~\footnote{\url{https://openai.com/blog/chatgpt-plugins\#\#code-interpreter}} for their superior capabilities, (2) Gemini Ultra and Pro \citep{gemini}, (3) Inflection-2 \citep{inflection-2}, (4) Grok-1~\footnote{\url{https://x.ai/model-card}}, and a group of models from Chinese firms, including (5) Baichuan-3~\footnote{\url{https://www.baichuan-ai.com}} and (6) the latest GLM-4~\footnote{\url{https://open.bigmodel.cn/dev/api\#glm-4}} from the GLM series \citep{glm}.
\end{itemize}

The models discussed here serve broad, general-purpose functions and have typically been subject to multiple alignment strategies. 
    \item \textbf{Open-source models} comprise: general-purpose systems such as (1) DeepSeek-LLM-Chat 67B \citep{deepseek-llm}, (2) Qwen 72B \citep{qwen}, (3) SeaLLM-v2 7B \citep{seallm}, and (4) ChatGLM3 6B \citep{chatglm3}; in addition, there are specialized models with mathematical capabilities, including (5) InternLM2-Math 20B~\footnote{\url{https://github.com/InternLM/InternLM-Math}}, an extension of InternLM2 with subsequent math-oriented training and instruction fine-tuning; (6) Math-Shepherd-Mistral 7B, which leverages PPO \citep{schulman2017proximal} on Mistral 7B \citep{mistral} utilizing a reward model grounded in process supervision; (7) the WizardMath family \citep{wizardmath}, which enhances Mistral 7B and Llama-2 70B \citep{llama2} in mathematical reasoning through a mix of evolve-instruct (instruction tuning driven by AI-generated instructions) and PPO with training items mostly sourced from GSM8K and MATH; (8) MetaMath 70B \citep{metamath}, achieved by fine-tuning Llama-2 70B on an enriched version of GSM8K and MATH; (9) ToRA 34B \cite{tora}, where CodeLlama 34B is adapted for tool-augmented mathematical reasoning; and (10) MAmmoTH 70B \citep{MathInstruct}, which applies instruction tuning based on MathInstruct to Llama-2 70B.
\end{itemize}

According to Table \ref{tab:sft_rl_math}, when tool-based problem solving is not permitted in evaluation, DeepSeekMath-Instruct 7B demonstrates notably strong stepwise reasoning skills. Specifically, on the challenging MATH benchmark, our model outperforms all publicly available open-source systems as well as the majority of proprietary solutions (including Inflection-2 and Gemini Pro), achieving a lead of at least 9\% in absolute accuracy. This holds true even in comparison to much larger models (like Qwen 72B) or those with dedicated reinforcement learning enhancements for mathematics (such as WizardMath-v1.1 7B). While DeepSeekMath-Instruct performs on par with Chinese proprietary models GLM-4 and Baichuan-3 on MATH, it still lags behind GPT-4 and Gemini Ultra.

In scenarios permitting a combination of natural language reasoning with programming tools for answering questions, DeepSeekMath-Instruct 7B attains nearly 60\% accuracy on MATH, exceeding all prior open-source entries. On additional evaluation sets, our model achieves results on par with DeepSeek-LLM-Chat 67B, which previously led performance among open-source models but is approximately ten times larger.

\section{Reinforcement Learning}

\subsection{Group Relative Policy Optimization} Reinforcement learning (RL) is well-established as an effective approach to further develop the mathematical problem-solving capabilities of large language models after initial Supervised Fine-Tuning (SFT) \citep{wang2023math,wizardmath}. Here, we present our novel RL algorithm, Group Relative Policy Optimization (GRPO), designed to be both resource-efficient and powerful.

\subsubsection{From PPO to GRPO} Proximal Policy Optimization (PPO) \citep{schulman2017proximal}, a widely-adopted actor-critic reinforcement learning algorithm, is frequently utilized for RL-based fine-tuning of large language models \citep{ouyang2022training}. PPO aims to maximize the following clipped surrogate objective:

\begin{equation}
\footnotesize
    \mathcal{J}_{PPO}(\theta) = \mathbb{E}{[q \sim P(Q), o \sim \pi_{\theta_{old}}(O|q)]} \frac{1}{|o|} \sum_{t=1}^{|o|} \min \left[ \frac{\pi_\theta(o_{t} | q, o_{<t})}{\pi_{\theta_{old}}(o_{t} | q, o_{<t})} A_{t}, \text{clip} \left( \frac{\pi_\theta(o_{t} | q, o_{<t})}{\pi_{\theta_{old}}(o_{t} | q, o_{

Here, $\pi_{\theta}$ and $\pi_{\theta_{old}}$ denote the current and prior policy networks, respectively, while $q$ and $o$ correspond to questions and responses drawn from the question corpus and sampled using the previous policy $\pi_{\theta_{old}}$. The parameter $\varepsilon$ serves as a clipping hyper-parameter, introduced in PPO to ensure training stability. $A_t$ represents the advantage, which is estimated using Generalized Advantage Estimation (GAE) \citep{gae} based on reward signals $\{r_{\ge t}\}$ and an associated learned value function $V_{\psi}$. Consequently, PPO necessitates joint training of both a value function and the policy network. To counteract excessive reward model optimization, the typical strategy is to augment the reward with a per-token KL divergence penalty relative to a reference model at every decoding step \citep{ouyang2022training}, specifically,

\begin{equation}
    r_{t} = r_\varphi(q, o_{\le t}) - \beta \log\frac{\pi_{\theta}(o_{t}|q, o_{<t})}{\pi_{ref}(o_{t}|q, o_{<t})},
\label{eq:PPO-reward}
\end{equation}

where $r_\varphi$ refers to the reward function, $\pi_{ref}$ is the reference policy (typically the initial SFT checkpoint), and $\beta$ regulates the KL divergence penalty.

Since the value function in PPO is generally realized by a separate network of similar scale as the policy model, this results in considerable computational and memory overhead. Moreover, during RL optimization, the value function serves as a baseline in the advantage computation to minimize variance. However, in large language model (LLM) scenarios, the reward model typically assigns a score only to the terminal token, posing challenges for training an accurate per-token value function. To overcome this, we introduce Group Relative Policy Optimization (GRPO), depicted in Figure \ref{fig:grpo}, which removes the requirement for a dedicated value function as used in PPO. Instead, GRPO takes the mean reward of several sampled completions, generated for the same prompt, as the baseline. Concretely, for each question $q$, GRPO samples a set of completions $\{o_1, o_2, \cdots, o_G\}$ from $\pi_{\theta_{old}}$, and maximizes the objective:

\begin{equation}
\footnotesize
\begin{split}
    \mathcal{J}_{GRPO}(\theta) &= \mathbb{E}{[q \sim P(Q), \{o_i\}_{i=1}^G \sim \pi_{\theta_{old}}(O|q)]}  \\
    & \frac{1}{G}\sum_{i=1}^G\frac{1}{|o_i|} \sum_{t=1}^{|o_i|} \left\{ \min \left[ \frac{\pi_\theta(o_{i,t} | q, o_{i,<t})}{\pi_{\theta_{old}}(o_{i,t} | q, o_{i,<t})} \hat{A}_{i,t}, \text{clip} \left( \frac{\pi_\theta(o_{i,t} | q, o_{i,<t})}{\pi_{\theta_{old}}(o_{i,t} | q, o_{i,<t})}, 1 - \varepsilon, 1 + \varepsilon \right)  \hat{A}_{i,t} \right] - \beta \mathbb{D}_{KL}\left[\pi_{\theta} || \pi_{ref}\right]\right\} ,
\end{split}
\label{eq:GRPO-obj}
\end{equation}

In this formulation, $\varepsilon$ and $\beta$ represent tunable hyper-parameters. The term $\hat{A}_{i,t}$ is the advantage, computed by comparing relative rewards only among outputs in the same group—a procedure detailed later. By leveraging grouped relative advantages, GRPO naturally fits the reward model's comparative training paradigm, which is typically built from datasets where outputs are rated relative to one another on the same prompt. Additionally, GRPO applies KL regularization directly to the loss term, penalizing the divergence between the learned and reference policies, rather than incorporating the penalty into the reward calculation—thereby streamlining the computation of $\hat{A}_{i,t}$. Distinct from the KL penalty adopted in (\ref{eq:PPO-reward}), we approximate the KL divergence using the following unbiased estimator, as per \citep{kl_approx}:

\begin{equation}
\small

\begin{algorithm}[t]
  \small
  \caption{Iterative Group Relative Policy Optimization}
  \textbf{Input} initial policy model $\pi_{\theta_{\text{init}}}$; reward models $r_\varphi$; task prompts $\mathcal{D}$; 
  hyperparameters $\varepsilon$, $\beta$, $\mu$
  \begin{algorithmic}[1]
    \State Initialize policy model $\pi_\theta \leftarrow \pi_{\theta_{\text{init}}}$
    \For{iteration = 1, \dots, I}
       \State Set reference model $\pi_{ref} \leftarrow \pi_{\theta}$
      \For{step = 1, \dots, M}
      \State Randomly select batch $\mathcal{D}_b$ from $\mathcal{D}$
      \State Update old policy model $\pi_{\theta_{old}} \leftarrow \pi_{\theta}$ 
      \State For each question $q \in \mathcal{D}_b$, generate $G$ samples $\{o_i\}_{i=1}^G \sim \pi_{\theta_{old}} (\cdot \mid q)$
      \State Compute corresponding rewards $\{r_i\}_{i=1}^{G}$ for the generated outputs $o_i$ using $r_{\varphi}$ 
      \State For every $t$-th token in $o_i$, evaluate $\hat{A}_{i,t}$ using group relative advantage estimation.
      \For{GRPO iteration = 1, \dots, $\mu$}
        \State Update the policy $\pi_{\theta}$ by optimizing the GRPO objective (Equation \ref{eq:GC-GRPO})
      \EndFor
    \EndFor \State Perform continual training of $r_\varphi$ leveraging a replay buffer.
    \EndFor 
  \end{algorithmic}
  \textbf{Output} $\pi_\theta$
  \label{alg:iter-grpo}
\end{algorithm}

\subsubsection{Outcome Supervision RL with GRPO} More precisely, given a question $q$, the prior policy $\pi_{\theta_{old}}$ is used to generate a collection of outputs $\{o_1, o_2, \cdots, o_G\}$. Each output is then assessed by the reward model, producing $G$ reward values $\mathbf{r}=\{r_1, r_2, \cdots, r_G\}$. These reward values are subsequently normalized: the mean of the group is subtracted and the result is divided by the group's standard deviation. Under outcome supervision, the terminal reward for each output $o_i$—now normalized—serves as the advantage for all tokens in that output, such that $\hat{A}_{i, t} = \widetilde{r}_i = \frac{r_i- {\rm mean}(\mathbf{r})}{{\rm std}(\mathbf{r})}$. The policy parameters are then refined by maximizing the objective in equation (\ref{eq:GRPO-obj}).

\subsubsection{Process Supervision RL with GRPO} Since outcome-based feedback provides a singular terminal reward, it might lack granularity or efficacy in complex mathematical problem-solving. Inspired by \cite{wang2023math}, this study also employs process supervision, delivering feedback after each reasoning stage. Concretely, given a question $q$ and the associated sampled outputs $\{o_1, o_2, \cdots, o_G\}$, a process-oriented reward model assigns a reward to each reasoning step, yielding the structured reward set $\mathbf{R} = \{ \{r_1^{index(1)},\cdots,r_1^{index(K_1)}\}, \cdots,  \{r_G^{index(1)},\cdots,r_G^{index(K_G)}\} \}$, where $index(j)$ identifies the terminal token of the $j$-th step and $K_i$ is the step count in the $i$-th output. Rewards at each step are normalized relative to the mean and standard deviation of all rewards: $\widetilde{r}_i^{index(j)} = \frac{r_i^{index(j)} - {\rm mean(\mathbf{R})}}{{\rm std(\mathbf{R})}}$. For each token, the process advantage is then the sum of normalized rewards for subsequent steps: $\hat{A}_{i, t} = \sum_{index(j) \ge t} \widetilde{r}_i^{index(j)}$. The optimization of the policy proceeds via maximizing the criterion in equation (\ref{eq:GRPO-obj}).

\subsubsection{Iterative RL with GRPO} During the reinforcement learning process, the previous reward model may eventually become inadequate to properly guide the current policy model. To address this, we also investigate an iterative RL framework using GRPO. As depicted in Algorithm \ref{alg:iter-grpo}, iterative GRPO works by continually creating new datasets for the reward model using outputs sampled from the evolving policy model. The old reward model is retrained by applying a replay mechanism, which retains 10\% of prior data. Subsequently, the policy model is used to update the reference model, after which further training is performed on the policy model using the updated reward model.

\subsection{Training and Evaluating DeepSeekMath-RL}

Our reinforcement learning experiments utilize DeepSeekMath-Instruct 7B as the policy model. The RL training data consists exclusively of chain-of-thought formatted items related to GSM8K and MATH that were sourced from the SFT data, totaling approximately 144K samples. We deliberately omit other SFT data types so as to isolate the effect of RL on evaluation benchmarks not included during the RL stage. The reward model training set is constructed according to \citep{wang2023math}. The initial reward model is trained from DeepSeekMath-Base 7B using a learning rate of 2e-5. For training the policy model with GRPO, we adopt a learning rate of 1e-6, and set the KL regularization coefficient at 0.04. For each question in the dataset, we sample $64$ candidate completions. Sequence length is capped at 1024 tokens, and the batch size for training is also 1024. After each round of exploration, the policy model undergoes just one update. We follow the DeepSeekMath-Instruct 7B evaluation methodology to assess DeepSeekMath-RL 7B. For DeepSeekMath-RL 7B, chain-of-thought reasoning tasks from GSM8K and MATH are classified as in-domain, while all other benchmarks are treated as out-of-domain tasks.

Table \ref{tab:sft_rl_math} presents a comparison of the performance of open- and closed-source models, evaluated with both chain-of-thought and tool-integrated reasoning approaches on English and Chinese benchmarks. The results indicate the following: 1) Using chain-of-thought reasoning, DeepSeekMath-RL 7B achieves accuracies of 88.2\% on GSM8K and 51.7\% on MATH. These outcomes exceed those of all open-source models within the 7B–70B parameter range and surpass most closed-source alternatives. 2) Notably, DeepSeekMath-RL 7B is exclusively trained on chain-of-thought formatted instruction tuning data from GSM8K and MATH, building on DeepSeekMath-Instruct 7B. Although its training set is limited in scope, DeepSeekMath-RL 7B consistently outperforms DeepSeekMath-Instruct 7B across every evaluation metric, demonstrating the efficacy of reinforcement learning.

\section{Discussion}
In the following section, we discuss key insights gained from our pre-training and RL experimental work.

\subsection{Insights from Pre-Training}

Here, we outline our experiences related to pre-training. Unless indicated otherwise, we maintain the training protocol detailed in Section \ref{sec:quality-policy}. It should be highlighted that references to the DeepSeekMath Corpus in this context pertain to the 89B-token dataset curated during the second round of data collection.

\subsubsection{Advantages of Code Training for Mathematical Reasoning}

There is a widely held but as yet insufficiently substantiated belief that exposure to code data enhances reasoning skills. We provide new evidence supporting this claim specifically in mathematical problem solving: code training strengthens models’ mathematical reasoning abilities both with and without auxiliary tool usage.

To evaluate the influence of code training on mathematical reasoning, we explored both a two-stage training approach and a one-stage training protocol as follows:

\textbf{Two-Stage Training}

\begin{itemize}[topsep=0pt]
    \item \textbf{Code Training for 400B Tokens $\rightarrow$ Math Training for 150B Tokens}: DeepSeek-LLM 1.3B undergoes initial training on 400B code tokens, followed by an additional 150B tokens consisting of mathematical content;
    \item \textbf{General Training for 400B Tokens $\rightarrow$ Math Training for 150B Tokens}: To provide a baseline, we perform a control experiment by substituting the 400B code tokens with 400B general-purpose tokens drawn from a large general text corpus curated by DeepSeek-AI. This enables an assessment of the specific contribution of code tokens relative to generic data for mathematical reasoning improvement.
\end{itemize}

\textbf{One-Stage Training}

\begin{itemize}[topsep=0pt]
    \item \textbf{Math Training for 150B Tokens}: We exclusively train DeepSeek-LLM 1.3B on a dataset of 150B mathematical tokens;
    \item \textbf{Training on a mixture of 400B Code Tokens and 150B Math Tokens}: Considering that consecutive math training after code training negatively affects coding ability, we examine a one-stage strategy where code and math tokens are interleaved. The purpose is to determine if this approach maintains gains in mathematical reasoning while reducing the extent of catastrophic forgetting.
\end{itemize}

\paragraph{Results}
Tables \ref{tab:code-math} and \ref{tab:code-math-general-eval} display model performance across the aforementioned training regimes.

Introducing code training enhances program-based mathematical reasoning regardless of whether training is staged or simultaneous. Specifically, Table \ref{tab:code-math} reveals that code pretraining in a two-stage approach yields notable gains in tackling GSM8K and MATH datasets using Python, which are then further increased by subsequent math-centric training. Notably, combining code and math tokens in a one-stage regimen substantially mitigates catastrophic forgetting, a challenge in two-stage protocols, while benefiting both general coding skills (Table \ref{tab:code-math-general-eval}) and mathematical reasoning reliant on programs (Table \ref{tab:code-math}).

Moreover, code pretraining shows value for pure mathematical reasoning tasks not involving external tools. In two-stage training, early code exposure confers moderate advantages and bolsters downstream math training effectiveness, leading to peak performance. In contrast, one-stage joint exposure to code and math appears to impair tool-free mathematical reasoning; we hypothesize that the limited capacity of DeepSeek-LLM 1.3B may prevent it from effectively integrating both domains when presented concurrently.

\subsubsection{ArXiv Papers Appear Unhelpful for Enhancing Mathematical Reasoning}

It is standard practice to incorporate ArXiv papers as part of the math pre-training corpus \citep{minerva,gpt-f,llemma,mathpile}. Yet, little comprehensive evaluation has been carried out to assess how their inclusion affects mathematical reasoning abilities. Contrary to prevailing assumptions, our empirical studies indicate that ArXiv papers do not significantly aid mathematical reasoning. Our investigation encompasses multiple model scales, such as DeepSeek-LLM 1.3B and DeepSeek-Coder-Base-v1.5 7B \citep{deepseek-coder}, and uses different processed forms of ArXiv corpora:

\begin{itemize}[topsep=0pt]
    \item \textbf{MathPile} \citep{mathpile}: an 8.9B-token dataset assembled via heuristic-based cleaning and filtering, where scientific ArXiv papers constitute more than 85\% of the collection;
    \item \textbf{ArXiv-RedPajama} \citep{redpajama}: comprising 28.0B tokens of complete ArXiv LaTeX content, stripped of preambles, commentary, macros, and bibliographies.
\end{itemize}

Our methodology entails independently pre-training DeepSeek-LLM 1.3B for 150B tokens and DeepSeek-Coder-Base-v1.5 7B for 40B tokens on each variant of the ArXiv corpus. Across the board, the use of ArXiv papers fails to yield noticeable gains—and in some instances leads to declines—in mathematical reasoning performance, as measured by a suite of benchmarks spanning various challenge levels. These assessments include quantitative datasets like GSM8K and MATH (see Table \ref{tab:arxiv-cot}), multiple-choice exams such as MMLU-STEM (Table \ref{tab:arxiv-cot}), and theorem-proving benchmarks like miniF2F (Table \ref{tab:arxiv_atp}).

Nevertheless, this assessment is not definitive and should be interpreted with caution. We have yet to explore:

\begin{itemize}[topsep=0pt]
    \item The influence of ArXiv-derived data on specialized math tasks absent from this investigation, for example, the informalization of theorems, which involves transforming formal proofs or statements into informal text;
    \item The outcomes arising from combining ArXiv data with alternative data types;
    \item Whether larger model capacities might reveal latent advantages from ArXiv-sourced pre-training.
\end{itemize}

Accordingly, there is ample scope for further inquiry, which we defer to subsequent research.

\subsection{Insights of Reinforcement Learning}
\subsubsection{Towards to a Unified Paradigm} This section presents a generalized paradigm for analyzing diverse training strategies—including SFT, RFT, DPO, PPO, and GRPO—and reports empirical studies examining elements within this framework. Typically, the gradient of the training objective with respect to parameters $\theta$ for a given method can be expressed as:

\begin{equation}
    \nabla_{\theta}\mathcal{J}_{\textcolor{red}{\mathcal{A}}}(\theta) = \mathbb{E}[\underbrace{(q,o) \sim \textcolor{red}{\mathcal{D}}}_{Data \ Source}]\left( \frac{1}{|o|} \sum_{t=1}^{|o|}  \underbrace{GC_{{\mathcal{A}}}(q, o, t, \textcolor{red}{\pi_{{rf}}})}_{Gradient \ Coefficient}  \nabla_{\theta}\log \pi_{\theta}(o_t | q, o_{<t})\right).
\label{eq:objective}
\end{equation}

This formulation centers around three main elements: 1) the \textit{Data Source $\mathcal{D}$}, specifying the collection from which training examples are drawn; 2) the \textit{Reward Function $\pi_{{rf}}$}, serving as the provider of feedback during learning; 3) the \textit{Algorithm $\mathcal{A}$}, which integrates the input data and the reward feedback into the gradient coefficient $GC$, thereby modulating the direction and scale of parameter updates. We categorize and interpret several common approaches within this unified conceptual framework:

\begin{itemize}
    \item  \textbf{Supervised Fine-tuning (SFT)}: In SFT, a pretrained model undergoes further training using manually curated SFT datasets.
    \item \textbf{Rejection Sampling Fine-tuning (RFT)}: RFT involves additional tuning of an SFT model by leveraging only those sampled outputs from the SFT model that pass correctness-based filtering criteria, using the same prompts as in SFT.
    \item \textbf{Direct Preference Optimization (DPO)}: DPO optimizes the SFT model by exposing it to a set of alternative responses from the SFT model and updating the model through a pairwise preference loss (DPO loss).
    \item \textbf{Online Rejection Sampling Fine-tuning (Online RFT)}: Unlike standard RFT, Online RFT uses the SFT-initialized policy as a starting point but refines the policy through ongoing fine-tuning with outputs generated and filtered in real time.
    \item \textbf{PPO/GRPO}: In PPO/GRPO, the policy network is initialized from the SFT model and further improved by applying reinforcement through samples generated by the current policy in an online fashion.
\end{itemize}

The components of these approaches are outlined in Table \ref{tab:data_policy}. For an expanded derivation, see Appendix \ref{app:analysis-rl}.

\paragraph{Reflection on Data Source} We classify data sources into two primary types: online and offline sampling. Online sampling refers to training data generated directly from the current training policy model during exploration, whereas offline sampling utilizes data gathered from the initial SFT model’s sampling outputs. Methods such as RFT and DPO adhere to offline sampling, while Online RFT and GRPO rely on online sampling approaches.

Figure \ref{fig:rl-analysis} presents evidence that Online RFT demonstrates marked performance improvements over RFT on both benchmarks. Notably, although Online RFT and RFT perform similarly at the outset of training, Online RFT achieves substantial gains in later stages, underscoring the strengths of online training regimes. This pattern is expected: early in training, the actor model remains similar to the SFT model, so their sampled data exhibit minimal divergence. As training progresses, the distinctions in the data sampled from the actor become more pronounced, and the advantage of sampling from the live actor model becomes more significant.

\paragraph{Reflection on Gradient Coefficient} The reward signal is converted to a gradient coefficient by the algorithm for updating model parameters. Our experiments distinguish between two types of reward functions: ‘Rule’ and ‘Model.’ The ‘Rule’ approach evaluates the correctness of responses based on predefined criteria, whereas the ‘Model’ approach involves training a reward model to assign scores to responses, using rule-based judgments as its training data. Equations \ref{eq:GC-RFT} and \ref{eq:GC-GRPO} reveal an essential distinction between GRPO and Online RFT: GRPO uniquely tailors its gradient coefficient in accordance with the reward magnitude received from the reward model, thereby applying variable reinforcement or penalty depending on the response’s evaluated quality. Conversely, Online RFT does not possess this adaptive mechanism; it applies consistent reinforcement for all correct answers and does not penalize incorrect ones, treating all correct responses uniformly.

As illustrated in Figure \ref{fig:rl-analysis}, GRPO consistently outperforms online RFT, underscoring the effectiveness of dynamically adjusting positive and negative gradient coefficients. Moreover, the combination GRPO+PS yields better results than GRPO+OS, suggesting that employing step-wise, granular gradient coefficients brings additional gains. We further investigate iterative RL, implementing two iterations in our studies. According to Figure \ref{fig:iter-rl}, iterative RL notably enhances performance, with the most pronounced improvement occurring after the first iteration.

\subsubsection{Why RL Works?} This study applies reinforcement learning to a selected subset of instruction tuning data, leading to substantial improvements over the baseline instruction-tuned model. To further analyze the impact of reinforcement learning, we examine both Pass@K and Maj@K accuracies for Instruct and RL models across two benchmark datasets. Figure \ref{fig:maj-pass} demonstrates that RL predominantly boosts Maj@K, while Pass@K remains largely unaffected. This observation suggests that RL primarily strengthens the robustness of the output distribution, i.e., \textbf{the observed gains mainly result from amplifying correct answers within the TopK rather than an enhancement of the model’s fundamental skill set.} Analogously, \citep{wang2023making} point out a \textbf{misalignment problem} in reasoning tasks handled by SFT models, and further show that SFT performance can benefit from diverse preference alignment techniques \citep{yuan2023rrhf,song2023preference,wang2023making}.

\subsubsection{How to Achieve More Effective RL?}
\label{sec:effec_rl}
Our experiments reveal that RL is highly effective for mathematical reasoning. Additionally, we introduce a cohesive framework to interpret a range of prominent training methodologies, recasting all these methods as either explicit or streamlined RL approaches. As distilled in Equation \ref{eq:objective}, this framework revolves around three critical elements: Data Source, Algorithm, and Reward Function. We outline several prospective directions for each component.

\paragraph{Data Source} The choice of data underpins all training approaches. Within the RL framework, the data source is defined as the set of unlabeled questions coupled with their policy-model-generated outputs. Our current work utilizes prompts from instruction tuning, relying on simple nucleus sampling to generate responses. We believe this limitation might explain why our RL approach selectively elevates Maj@K performance. Going forward, we plan to extend our RL framework to encompass out-of-distribution queries, as well as leverage \textbf{advanced sampling (decoding) strategies}—such as those employing tree-search algorithms \citep{yao2023tree}. In addition, \textbf{efficient inference techniques} \citep{xia-etal-2023-speculative,leviathan2023fast,kwon2023efficient,xia2024unlocking}—which govern how efficiently the policy model explores its environment—are recognized as crucial to further performance gains.

\paragraph{Algorithms} Algorithms utilize both the data and the reward signals to determine the gradient coefficients necessary for updating model parameters. According to Equation \ref{eq:objective}, contemporary techniques predominantly \textbf{TRUST} the reward function's signal, adjusting the conditional probability of specific tokens upwards or downwards accordingly. Nonetheless, the assumption that the reward signal is invariably accurate does not always hold true, particularly for tasks characterized by high complexity. As an illustration, the PRM800K datasets \citep{lightman2023let}, even though meticulously labeled by skilled annotators, still exhibit about 20\% annotation errors\footnote{\url{https://github.com/openai/prm800k/issues/12\#issuecomment-1728491852}}. Therefore, we aim to investigate reinforcement learning algorithms that demonstrate resilience in the presence of noisy reward signals. Our perspective is that \textbf{WEAK-TO-STRONG} \citep{burns2023weak} alignment techniques have the potential to fundamentally transform learning algorithm paradigms.

\paragraph{Reward Function} The reward function delivers the principal signal for model training. Within RL frameworks, this function is most often instantiated as a neural reward model. We identify three pivotal research avenues for advancing reward models: 1) \textbf{Enhancing the generalization capacity of reward models.} It is imperative that reward models can generalize to unseen data distributions and cope with sophisticated decoding outputs; without this, reinforcement learning risks merely preserving the existing distribution of LLMs rather than enabling substantial improvements; 2) \textbf{Quantifying the uncertainty inherent in reward models.} Uncertainty metrics could serve as a crucial link, connecting imperfect reward models with weak-to-strong algorithmic strategies; 3) \textbf{Efficient construction of high-fidelity process reward models} that deliver detailed training signals during the step-by-step reasoning process \citep{lightman2023let,wang2023math}.

\section{Conclusion, Limitation, and Future Work} This work introduces DeepSeekMath, which surpasses all currently available open-source systems on the MATH benchmark designed for competition-level problems and achieves results comparable to those of proprietary models. The initialization of DeepSeekMath~is based on DeepSeek-Coder-v1.5 7B, followed by continual training on a corpus totaling 500B tokens, where math-specific content from Common Crawl contributes 120B tokens—a substantial fraction. Through detailed ablation experiments, we demonstrate that websites are a promising reservoir for obtaining high-caliber mathematical content; meanwhile, the anticipated advantage of arXiv as a data source does not materialize to the same extent. We also present Group Relative Policy Optimization (GRPO), an adaptation of Proximal Policy Optimization (PPO), which enables marked improvements in mathematical reasoning while reducing memory overhead. Our empirical analysis confirms GRPO’s efficacy, even when DeepSeekMath-Instruct 7B achieves already strong benchmark results. Furthermore, we propose a comprehensive framework to conceptualize various related methodologies, and identify multiple promising directions to further refine reinforcement learning effectiveness.

Despite the notable achievements of DeepSeekMath~in tasks involving quantitative reasoning, its performance on geometry and theorem-proof remains inferior relative to closed-source counterparts. For example, preliminary experiments show the model struggles with triangle- and ellipse-related questions, potentially due to biases introduced in both pre-training and fine-tuning data curation. Moreover, owing to its parameter size, DeepSeekMath~lags behind GPT-4 in few-shot settings; while GPT-4 demonstrates notable gains with few-shot prompts, DeepSeekMath~displays minimal variation between zero-shot and few-shot evaluations. Going forward, our focus will be on refining the engineered data selection process to build a more robust and high-quality pre-training dataset. We also intend to pursue the research directions outlined in Section \ref{sec:effec_rl} to further advance reinforcement learning for large language models.

\end{document}